\section{Positive frame weight and the Euler direction}
\label{bl:sec:frames}

We now combine the graph module with the scaling action on the
frame bundle. A positive weight has two consequences: after a
transverse cover it gives an ample source for the horizontal symbol,
and in the vertical direction it gives a nonzero Euler eigenvalue.
The proof keeps these two operations separate until both have been
identified on the same right differential module.

Return to Section~\ref{bl:sec:covers}. Thus $T=\PP^\ell$,
$J=\OO_T(1)$, and $B=J^\times$ is the bundle of
nonzero vectors, with projection $p_B:B\to T$. The parameter
$\lambda\in\Gm$ acts on $B$ by multiplication by $\lambda^r$.
Its action on $Y$ preserves $R$ and $D_h$, and $h:R\to B$ is
equivariant. All the modules in~\eqref{bl:hod:modules} inherit their
natural geometric action.

Let $\varepsilon$ be the infinitesimal generator on $B$. In a
bundle chart with fiber coordinate $\beta$,
\begin{equation}\label{bl:frm:euler}
 \varepsilon=r\beta\frac{\partial}{\partial\beta}.
\end{equation}
It is nowhere zero and spans the kernel of
$T_B\to p_B^*T_T$. By weight $m$ we mean pullback weight:
$\lambda^*s=\lambda^m s$.

\begin{proposition}[Positive-weight obstruction]
\label{bl:prop:frame-positivity}
For the filtered module $\mathcal N$ of~\eqref{bl:hod:modules}, the
following assertions hold.
\begin{enumerate}[label=\textup{(\roman*)}]
\item If
$v\in H^0(B,F_0\mathcal N\otimes T_B)$ is homogeneous of weight
$m>0$ and has zero symbol in $\gr_1^F\mathcal N$, then its image
in $F_0\mathcal N\otimes p_B^*T_T$ is zero.
\item If $s\in H^0(B,F_0\mathcal N)$ is homogeneous of weight $m$,
then its actual right differential-operator action satisfies
\begin{equation}\label{bl:frm:euler-action}
 s\varepsilon=-ms.
\end{equation}
The same assertion is valid on invariant bundle charts.
\end{enumerate}
\end{proposition}

The first assertion concerns the symbol, and the second concerns the
operator itself. Their combination will annihilate the obstruction.
We prove them separately, keeping track of both canonical factors
and filtration shifts.

The horizontal directions cannot in general be suppressed. For
example, if $V=\PP^2$ and $C$ is a smooth quartic, then
$A=K_V+C\sim\OO_{\PP^2}(1)$ and $G=C\sim4A$ satisfy the
supported hypotheses, while $\OO_V(G)|_C$ is ample and nontrivial.
The construction must therefore retain the full map $f:S\to T$.

\subsection{A transverse power cover}

Suppose first that $\ell>0$. Take the coordinate $r$-th power map of
$\PP^\ell$ and compose it with a general target automorphism, obtaining
\[
 \phi:T'=\PP^\ell\longrightarrow T,
 \qquad J'=\OO_{T'}(1),\qquad (J')^r\simeq\phi^*J.
\]
It is finite flat. Put $B'=(J')^\times$, with its ordinary scalar
action, and let $p':B'\to T'$ be the projection. Raising frames to
the $r$-th power gives a finite flat equivariant map $\eta$ in
\begin{equation}\label{bl:frm:bundle-square}
\begin{tikzcd}[column sep=large,row sep=large]
B'=(J')^\times \arrow[r,"\eta"] \arrow[d,"p'"'] &
B=J^\times \arrow[d,"p_B"]\\
T' \arrow[r,"\phi"'] & T.
\end{tikzcd}
\end{equation}
Along each nonzero fiber its differential is invertible.

Choose $\phi$ transverse to the maps to $T$ from every closed
intersection of components of $R+D_h$ involving a component of $R$.
These intersections are smooth and there are finitely many of them.
Here is the generic-transversality justification. For such an
intersection $L$, the matching-pair space for the map $L\to T$
and the translated power map is smooth: variation of the target
automorphism supplies every target tangent direction. A general
fiber over the automorphism group is smooth by generic smoothness
in characteristic zero. This is exactly the required transversality.
The finitely many conditions can be imposed simultaneously.
Since the vertical differential in~\eqref{bl:frm:bundle-square} is
invertible, this also gives transversality to the corresponding
maps $L\to B$. If a point of $T$ is prescribed, one may in addition
require that the branch divisor avoid it, an open condition on
the target automorphism.

Let $Q'=Q\times_B B'\subset Y\times B'$, and set
\begin{equation}\label{bl:frm:basechanged-modules}
 \mathcal M'=\mathcal H^{b+1}_{[Q']}
     \bigl(\omega_{Y\times B'}(*D_h)\bigr),
 \qquad
 \mathcal N'=\mathcal H^1\operatorname{pr}_{B',+}\mathcal M'.
\end{equation}
These are again mixed Hodge modules with projective support.

\begin{lemma}[Filtered transverse comparison]\label{bl:frm:comparison}
For every $j$ there are natural equivariant isomorphisms
\begin{equation}\label{bl:frm:filtered-comparison}
 F_j\mathcal N'
   \simeq \eta^*F_j\mathcal N\otimes\omega_{B'/B}.
\end{equation}
Under these isomorphisms, symbol action transforms by the tangent
map $d\eta:T_{B'}\to\eta^*T_B$.
\end{lemma}

\begin{proof}
Use the local coordinates in the proof of
Lemma~\ref{bl:lem:lowest-piece}. The pulled-back graph equations are
\[
 \psi'_i=t_i\circ\eta-v_i.
\]
At a point of $Q'$, transversality to the intersection of all
boundary components through that point says that the differentials
of the $\psi'_i$, together with those boundary coordinates, are
independent. They can therefore be used as part of a local smooth
coordinate system. In particular the local support calculation
remains a normal crossing calculation.

Flat pullback preserves the localization quotient and its local
cohomology. More specifically, the explicit description following
Lemma~\ref{bl:lem:lowest-piece} identifies its filtration level by
level: the fractions with total extra denominator exponent at most
$j$ pull back to precisely the corresponding fractions for $\mathcal M'$.
For right modules one must also replace the absolute top-form
frame on $Y\times B$ by that on $Y\times B'$. Thus the filtered
comparison before direct image is
\[
 F_j\mathcal M'
  \simeq (\id_Y\times\eta)^*F_j\mathcal M
                \otimes\operatorname{pr}_{B'}^*\omega_{B'/B}.
\]
This is an isomorphism with the abstract relative canonical line;
it does not multiply sections by a Jacobian that vanishes at
ramification.

The relative de Rham differentials differentiate only in $Y$.
They commute with this comparison and with its relative canonical
factor. Flat base change therefore applies to the entire filtered
relative de Rham complex. One can check it on a finite affine
cover and its \v{C}ech complex; flat pullback commutes with all
cohomology in question. Projective-support strictness identifies
the cohomology filtration with that of $\mathcal N$ and
$\mathcal N'$, giving~\eqref{bl:frm:filtered-comparison} for all $j$.
The chain rule in the parameter coordinates gives the assertion
about symbols. Terms arising from a change of canonical frame
have order zero and do not change that assertion.
\end{proof}

One can also express the transversality in this proof as a
non-characteristic condition. The characteristic variety of the
normal crossing graph module is contained in the union of
conormals to the graph strata involving $R$. A characteristic
covector killed by $d(\id_Y\times\eta)$ has zero $Y$ component;
its remaining component annihilates both the stratum tangent
image and $d\eta(T_{B'})$. Transversality forces it to be zero.
Finite flatness also makes the higher $\operatorname{Tor}$ groups of
every graded filtration piece vanish. These are the two
non-characteristic conditions of
\cite[Section~3.5.1 and Lemmas~3.5.3, 3.5.5]{Saito1988}; in relative dimension
zero they give the same filtration index and right canonical factor
as the explicit calculation above.

\subsection{The adjugate and the projective quotient}

For vector bundles of equal rank the adjugate of $d\eta$ defines
a regular morphism
\begin{equation}\label{bl:frm:adjugate}
 \operatorname{adj}(d\eta):
 \eta^*T_B\longrightarrow\omega_{B'/B}\otimes T_{B'}.
\end{equation}
Indeed $\det(d\eta)$ is a section of $\omega_{B'/B}$, so this
is the usual determinant-twisted adjugate. It involves no division.
If $v$ has zero symbol, applying~\eqref{bl:frm:adjugate} to its tangent
factor and then~\eqref{bl:frm:filtered-comparison} gives a section
\begin{equation}\label{bl:frm:transported-symbol}
 v'\in H^0(B',F_0\mathcal N'\otimes T_{B'})
\end{equation}
with zero symbol. In local matrices this is just
$d\eta\operatorname{adj}(d\eta)=\det(d\eta)\id$.
The construction preserves the weight of $v$.

Suppose the horizontal image of $v$ is nonzero. Choose a point
where that coherent-sheaf section is nonzero, and choose $\phi$
as above with branch divisor avoiding its image in $T$.
On an unbranched neighborhood the tangent map is an isomorphism
preserving the vertical tangent line. Faithful flatness there
shows that the horizontal image of $v'$ is still nonzero.

Now take the diagonal quotient
\[
 \overline X=(Y\times B')/\Gm\longrightarrow T'.
\]
It exists as an associated algebraic variety: a local section of
$B'\to T'$ identifies it with $Y$ times that base chart, and changes
of section glue these products through the action on $Y$.
It is smooth. The linearized ample bundle on $Y$ descends to a
relatively ample bundle on this associated variety, so
$\overline X$ is quasi-projective. The invariant divisor $D_h$
descends to a divisor $\overline D_h$, and the invariant closed
support $Q'$ descends to $\overline Q$.

The support map $\overline Q\to T'$ is projective. Properness
descends from the projective map $Q'\to B'$, and a descended
relatively ample bundle gives projectivity. The quotient changes
both the ambient dimension and the support dimension by one;
thus $\overline Q$ still has codimension $b+1$. Define
\begin{equation}\label{bl:frm:quotient-modules}
 \overline{\mathcal M}=
 \mathcal H^{b+1}_{[\overline Q]}
       \bigl(\omega_{\overline X}(*\overline D_h)\bigr),
 \qquad
 \overline{\mathcal N}
       =\mathcal H^1\overline{\operatorname{pr}}_+
                          \overline{\mathcal M}.
\end{equation}
Use the constant twist by $\dim\overline X$ in this formula.
These are algebraic mixed Hodge modules, and the second has
projective support over the projective base $T'$.

Choose a frame of $J'$ on $U\subset T'$ and write a point of
$Y\times B'|_U$ as $(z,t,u)$, with $z\in Y$, $t\in U$, and
$u\in\Gm$. The coordinate change
\begin{equation}\label{bl:frm:untwisting}
 (z,t,u)\longmapsto(u^{-1}\!\cdot z,t,u)
\end{equation}
turns the diagonal action
$(z,t,u)\mapsto(\lambda\!\cdot z,t,\lambda u)$ into the action
on $u$ alone. Since $Q'$ and $D_h$ are invariant, the support
and divisor become products with $\Gm$ as well. Localization
and support commute with this product. Their normalizations give
\begin{equation}\label{bl:frm:quotient-filtration}
 F_j\mathcal N'
   \simeq (p')^*F_j\overline{\mathcal N}
                     \otimes\omega_{B'/T'}.
\end{equation}
To check the index, ordinary right smooth pullback of relative
dimension one changes $F_j$ to $F_{j+1}$ and tensors with the
relative canonical line. The normalized constant upstairs has
one additional Tate twist. This cancels the index change: the
normalized smooth inverse image is $[1](1)$, and gives exactly
\eqref{bl:frm:quotient-filtration}. See~\cite[Section~30]{Schnell2014}.
Relative direct image has the same relative dimension before and
after this product, so its cohomological index stays one.

The relative canonical line in~\eqref{bl:frm:quotient-filtration}
has invariant frame $du/u$. This frame is intrinsic as a relative
form: a change of bundle trivialization changes it only by a
form from the base. In the product description the vertical
symbol is zero and the horizontal symbol is the pullback of that
on $\overline{\mathcal N}$. In particular
$F_{<0}\overline{\mathcal N}=0$, by~\eqref{bl:frm:filtered-comparison},
\eqref{bl:frm:quotient-filtration} and faithful flatness.

Project $v'$ horizontally and use~\eqref{bl:frm:quotient-filtration}.
A homogeneous section of weight $m$ on the nonzero vectors of
$J'$ corresponds to a section twisted by $(J')^{-m}$. Explicitly,
if local frames satisfy $e_j=g_{ij}e_i$ and a vector is $u_ie_i$,
then $u_j=g_{ij}^{-1}u_i$; equality of expressions $u_i^m s_i$
forces $s_j=g_{ij}^m s_i$, the transition rule for a homomorphism
from $(J')^m$. Thus the projected zero-symbol section gives
\begin{equation}\label{bl:frm:ample-map}
 (J')^m\longrightarrow
 \ker\left[F_0\overline{\mathcal N}\otimes T_{T'}
                   \longrightarrow\gr_1^F\overline{\mathcal N}\right].
\end{equation}
It is nonzero if the original horizontal projection was nonzero.
For $m>0$ its source is ample, contradicting
Lemma~\ref{bl:lem:symbol-vanishing}. This proves assertion
\textup{(i)} of Proposition~\ref{bl:prop:frame-positivity}.
For $\ell=0$ there is no horizontal tangent direction and
assertion~\textup{(i)} is immediate.

\subsection{The actual Euler action}

It remains to prove assertion~\textup{(ii)}, which does not follow
merely from the vanishing of the vertical symbol. Work on a
bundle chart of $B$ and make the finite \emph{\'etale} fiber cover
\[
 \beta=u^r,\qquad u\ne0.
\]
The generator~\eqref{bl:frm:euler} becomes $u\partial_u$.
The coordinate change~\eqref{bl:frm:untwisting} again identifies the
supported geometric construction with a product in the $u$ direction. This is an
identification of the underlying differential modules, with their
natural equivariance, not just an identification of coherent
filtration pieces.

A weight-$m$ section is consequently $u^m$ times a section
independent of $u$, in the invariant relative frame $du/u$.
Right action on top forms is negative Lie differentiation.
Since $du/u$ is invariant,
\[
 (u^m s_1\otimes du/u)\,(u\partial_u)
       =-m(u^m s_1\otimes du/u).
\]
The product direct image retains this action, proving
\eqref{bl:frm:euler-action} on the \'{e}tale cover and hence downstairs.
This completes the proof of Proposition~\ref{bl:prop:frame-positivity}.

\begin{remark}\label{bl:frm:operator-order}
If $v=s_0\otimes\varepsilon$ and
$\varepsilon=\sum_i c_i\partial_{t_i}$, then the relevant
right operator is
\[
 \sum_i(s_0c_i)\partial_{t_i}=s_0\varepsilon.
\]
The coefficients precede the derivatives. Moving them to the
other side would introduce a divergence correction; the displayed
identity is simply associativity in the ring of differential
operators.
\end{remark}
